\subsection{Transient effects}

\label{sec:transient}
 
\subsubsection*{Larmor Precession in the Rotating Frame}
 
In the Zeeman region the resonance frequency is :
\begin{equation}
    \omega_0 = \gamma B_0, \qquad \gamma \equiv \frac{g_F\,\mu_0}{\hbar}
    \label{eq:larmor}
\end{equation}
where $\gamma$ is the gyromagnetic ratio. The equation of motion of the total angular
momentum $\mathbf{F}$ in the presence of both a static field $B_0\hat{z}$ and an RF
field $\mathbf{B}_{\text{rf}}$ is :
\begin{equation}
    \frac{d\mathbf{F}}{dt} = \gamma\,\mathbf{F} \times \mathbf{B}
    \label{eq:eom}
\end{equation}
Transforming to a frame rotating at angular frequency $\omega$ about $\hat{z}$, the
effective field becomes :
\begin{equation}
    \mathbf{B}_{\text{eff}} = \mathbf{B}_0 + \frac{\boldsymbol{\omega}}{\gamma}
    = \left(B_0 - \frac{\omega}{\gamma}\right)\hat{z} + B_{\text{rf}}\hat{x}'
    \label{eq:Beff}
\end{equation}
so that the equation of motion in the rotating frame reduces to :
\begin{equation}
    \frac{\partial \mathbf{F}}{\partial t} = \gamma\,\mathbf{F} \times \mathbf{B}_{\text{eff}}
    \label{eq:eom_rot}
\end{equation}
The magnitude of the effective field is :
\begin{equation}
    B_{\text{eff}} = \sqrt{\left(B_0 - \frac{\omega}{\gamma}\right)^2 + B_{\text{rf}}^2}
    \label{eq:Beff_mag}
\end{equation}
with the tipping angle :
\begin{equation}
    \cos\theta = \frac{\omega_0 - \omega}{\gamma B_{\text{eff}}},
    \qquad
    \sin\theta = \frac{\gamma B_{\text{rf}}}{\gamma B_{\text{eff}}} = \frac{B_{\text{rf}}}{B_{\text{eff}}}
    \label{eq:theta}
\end{equation}
At exact resonance $\omega = \omega_0$, we have $\theta = 90°$ and
$\mathbf{B}_{\text{eff}} = B_{\text{rf}}\hat{x}'$, so $\mathbf{F}$ precesses about
$B_{\text{rf}}$ at the Rabi frequency :
\begin{equation}
    \Omega_R = \gamma B_{\text{rf}}
    \label{eq:rabi_freq}
\end{equation}
 
\subsubsection*{Rate Equations for Two-Level System}
 
Considering only the $M = +2$ and $M = +1$ sublevels, let $p_2$ and $p_1 = 1 - p_2$
be their occupation probabilities. With equal up/down transition probabilities $b_{12} = b_{21} \equiv b$ :
\begin{align}
    \dot{p}_2 &= -b\,p_2 + b\,p_1 = -b(2p_2 - 1) \label{eq:rate2} \\
    \dot{p}_1 &= +b\,p_2 - b\,p_1 = +b(2p_2 - 1) \label{eq:rate1}
\end{align}
Substituting $q \equiv 2p_2 - 1$ :
\begin{equation}
    \dot{q} = -2b\,q
    \label{eq:q_ode}
\end{equation}
with the solution :
\begin{equation}
    q(t) = q_0\,e^{-2bt}
    \quad\Longrightarrow\quad
    p_2(t) = \frac{1}{2} + \delta\,e^{-2bt}
    \label{eq:p2_solution}
\end{equation}
where $\delta = p_2(0) - \tfrac{1}{2}$ is the initial excess population. The RF thus
drives $p_2$ exponentially toward $\tfrac{1}{2}$, with the $1/e$ time :
\begin{equation}
    t_{1/e} = \frac{1}{2b}
    \label{eq:t1e}
\end{equation}
Since $b \propto B_{\text{rf}}^2$, the decay time is \emph{inversely proportional} to
the RF power.
 
\subsubsection*{Optical Pumping Transient (RF Off $\to$ On)}
 
Before the RF is applied, optical pumping has established :
\begin{equation}
    p_2(0) = \frac{1}{2} + \delta_0, \qquad p_1(0) = \frac{1}{2} - \delta_0
    \label{eq:initial_pop}
\end{equation}
When the RF is suddenly switched on at resonance, the net downward transition rate
exceeds the upward rate by :
\begin{equation}
    \dot{p}_2\big|_{t=0^+} = -2b\,\delta_0 < 0
    \label{eq:initial_rate}
\end{equation}
causing an immediate decrease in transmitted light intensity.
 
\subsubsection*{Rabi Oscillations (Damped Ringing)}
 
In the coherent limit (before relaxation damps the motion), $\mathbf{F}$ precesses about
$\mathbf{B}_{\text{eff}}$ at the Rabi frequency $\Omega_R$. The time-dependent
population is :
\begin{equation}
    p_2(t) = \frac{1}{2} + \left(\frac{1}{2} - p_2^{\text{eq}}\right)\cos(\Omega_R t)\,e^{-t/T_2}
    \label{eq:rabi_osc}
\end{equation}
where $T_2$ is the transverse relaxation time. The period of the ringing is therefore :
\begin{equation}
    T_{\text{ring}} = \frac{2\pi}{\Omega_R} = \frac{2\pi}{\gamma B_{\text{rf}}}
    \label{eq:ring_period}
\end{equation}
which is \emph{inversely proportional} to $B_{\text{rf}}$ (and hence to the RF current
in the coils). For the two isotopes at the same $B_{\text{rf}}$:
\begin{equation}
    \frac{T_{85}}{T_{87}} = \frac{\gamma_{87}}{\gamma_{85}}
    = \frac{g_F(^{87}\text{Rb})}{g_F(^{85}\text{Rb})}
    = \frac{1/2}{1/3} = \frac{3}{2}
    \label{eq:period_ratio}
\end{equation}
 
\subsubsection*{Optical Pumping Time (RF On $\to$ Off)}
 
When the RF is switched off, optical pumping rebuilds the population imbalance. The
transmitted intensity recovers exponentially :
\begin{equation}
    I(t) = I_{\max}\left(1 - e^{-t/\tau_p}\right) + I_{\min}\,e^{-t/\tau_p}
    \label{eq:pump_recovery}
\end{equation}
where the optical pumping time $\tau_p$ satisfies :
\begin{equation}
    \frac{1}{\tau_p} = \sum_j b_{j,\text{pump}} + \Gamma_{\text{rel}}
    \label{eq:pump_time}
\end{equation}
with $b_{j,\text{pump}}$ the photon absorption rate and $\Gamma_{\text{rel}}$ the
collisional relaxation rate. Since $b_{j,\text{pump}} \propto I_{\text{light}}$, the
pumping time $\tau_p$ is \emph{inversely proportional} to the pumping light intensity.
Experimentally, $\tau_p$ is of the order of 10--20~ms.